\chapter{Pendiferensialan}\label{ch:b1:derivative}

\omterm{def:b1:derivative:def}{Turunan} sudah dihitung di sepanjang jilid Sekolah Menengah; yang
kurang adalah rantai teorema yang mengubah perhitungan menjadi
keterangan tentang fungsinya: yaitu teorema Rolle, teorema nilai rata-rata,
dan akibatnya --- kriteria kemonotonan, batas Lipschitz, dan
kecembungan. Segalanya dalam bab ini menyangkut fungsi yang terdefinisi pada
sebuah \omterm{prop:b1:reals:intervals}{selang} $I$.

\section{Turunan}

\begin{definition}\label{def:b1:derivative:def}
Fungsi $f \colon I \to \R$ disebut \emph{dapat diturunkan di $x_0 \in
I$}\index{turunan} bila hasil bagi selisihnya
$\frac{f(x) - f(x_0)}{x - x_0}$ mempunyai limit (yang hingga) ketika $x \to x_0$;
dan limitnya ditulis $f'(x_0)$. Setara dengan itu:
\[
f(x_0 + h) = f(x_0) + f'(x_0)\,h + h\,\varepsilon(h),
\qquad \varepsilon(h) \xrightarrow[h \to 0]{} 0 ,
\]
sehingga grafiknya lalu mempunyai garis singgung $y = f(x_0) + f'(x_0)(x -
x_0)$. Sifat dapat diturunkan di $x_0$ mengakibatkan \omterm{def:b1:continuity:continuous}{kekontinuan} di $x_0$ (bacalah
tampilannya). Adapun $f$ dapat diturunkan pada $I$ bila ia demikian di setiap titiknya;
lalu $f$ berkelas $C^1$ bila lebih lanjut $f'$ \omterm{def:b1:continuity:continuous}{kontinu}, dan berkelas
$C^k$\index{fungsi Ck@fungsi $C^k$} bila $f$ dapat
diturunkan $k$ kali dengan $f^{(k)}$ yang \omterm{def:b1:continuity:continuous}{kontinu}.
\end{definition}

\begin{example}\label{ex:b1:derivative:notc1}
Konvers ``\omterm{def:b1:derivative:def}{dapat diturunkan} $\Rightarrow$ \omterm{def:b1:continuity:continuous}{kontinu}'' gagal:
$\abs{\,\cdot\,}$ di $0$. Yang lebih mengejutkan, \omterm{def:b1:derivative:def}{dapat diturunkan} tak
mengakibatkan $C^1$: karena fungsi $f(x) = x^2 \sin\frac 1x$ ($f(0) = 0$) bersifat
\omterm{def:b1:derivative:def}{dapat diturunkan} di mana-mana, dengan $f'(0) = 0$, tetapi $f'(x) = 2x
\sin\frac1x - \cos\frac 1x$ tak berlimit di $0$
(\cref{exo:b1:derivative:2}).
\end{example}

\begin{example}[Dapat diturunkan tepat di satu titik]\label{ex:b1:derivative:onepoint}
Misalkan $f(x) = x^2$ untuk $x \in \Q$ dan $f(x) = 0$ untuk $x \notin
\Q$. Di $0$: $\bigl|\frac{f(h) - 0}{h}\bigr| \leq \abs h \to 0$,
sehingga $f$ \omterm{def:b1:derivative:def}{dapat diturunkan} di $0$ dengan $f'(0) = 0$. Adapun di sebarang $x_0
\neq 0$, $f$ bahkan tak \omterm{def:b1:continuity:continuous}{kontinu}: karena barisan rasional dan irasional
yang konvergen ke $x_0$ mengirim $f$ ke $x_0^2 \neq 0$ dan ke
$0$ berturut-turut (menurut \omterm{def:b1:topology:dense}{kepadatannya}, \cref{thm:b1:reals:density}). Jadi
sifat \omterm{def:b1:derivative:def}{dapat diturunkan} sungguh gagasan yang \emph{titik demi titik}: ia dapat
berlaku di satu titik $\R$ dan tak di mana pun lagi. Adapun moralnya bagi
praktik: bahwa \omterm{def:b1:logic:statement}{pernyataan} seperti kriteria kemonotonan atau Rolle
menuntut \omterm{def:b1:derivative:def}{turunannya} \emph{pada sebuah \omterm{prop:b1:reals:intervals}{selang}} --- karena memiliki
$f'(x_0)$ di titik yang terasing, sebanyak apa pun, tak menyangga kesimpulan
menyeluruh apa pun.
\end{example}

\begin{theorem}[Operasi]\label{thm:b1:derivative:operations}
Jika $f, g$ \omterm{def:b1:derivative:def}{dapat diturunkan} di $x_0$ (dan di tempat rumusnya
bermakna):
\[
(f + g)' = f' + g', \qquad
(fg)' = f'g + fg', \qquad
\Bigl(\frac fg\Bigr)' = \frac{f'g - fg'}{g^2},
\]
lalu jika $g$ \omterm{def:b1:derivative:def}{dapat diturunkan} di $f(x_0)$:
$\;(g \circ f)'(x_0) = g'\bigl(f(x_0)\bigr)\, f'(x_0)$ (yaitu aturan
rantai\index{aturan rantai}).
\end{theorem}

\begin{proof}
Jumlahnya: langsung. Hasil kalinya: tulislah
\[
f(x)g(x) - f(x_0)g(x_0)
= \bigl(f(x) - f(x_0)\bigr) g(x) + f(x_0)\bigl(g(x) - g(x_0)\bigr),
\]
lalu bagilah dengan $x - x_0$ dan biarkan $x \to x_0$ ($g$ \omterm{def:b1:continuity:continuous}{kontinu} di $x_0$).
Hasil baginya: tanganilah $\frac 1g$ lewat $\frac{1/g(x) - 1/g(x_0)}{x - x_0} =
\frac{-1}{g(x)g(x_0)}\cdot\frac{g(x) - g(x_0)}{x - x_0}$, lalu terapkan
aturan hasil kalinya. Aturan rantainya: dengan $y_0 = f(x_0)$, definisikan
$\theta(y) = \frac{g(y) - g(y_0)}{y - y_0}$ untuk $y \neq y_0$ dan
$\theta(y_0) = g'(y_0)$: maka $\theta$ \omterm{def:b1:continuity:continuous}{kontinu} di $y_0$, dan untuk
$x \neq x_0$,
\[
\frac{g(f(x)) - g(f(x_0))}{x - x_0}
= \theta\bigl(f(x)\bigr)\cdot \frac{f(x) - f(x_0)}{x - x_0}
\longrightarrow g'(y_0)\, f'(x_0),
\]
dengan faktor pertamanya lewat komposisi limitnya (dan muslihat ini menangani
kasus $f(x) = f(x_0)$ dengan bersih, yang di situ ``kalikan lalu bagilah
dengan $f(x) - f(x_0)$'' yang naif itu patah).
\end{proof}

\begin{theorem}[Turunan fungsi invers]\label{thm:b1:derivative:inverse}
Misalkan $f$ \omterm{def:b1:continuity:continuous}{kontinu} dan monoton tegas pada $I$, \omterm{def:b1:derivative:def}{dapat diturunkan}
di $x_0$ dengan $f'(x_0) \neq 0$. Maka $f^{-1}$
(\cref{thm:b1:continuity:bijection}) \omterm{def:b1:derivative:def}{dapat diturunkan} di $y_0 =
f(x_0)$, dengan
\[
(f^{-1})'(y_0) = \frac{1}{f'(x_0)} = \frac{1}{f'\bigl(f^{-1}(y_0)\bigr)} .
\]
Sedangkan jika $f'(x_0) = 0$, maka inversnya mempunyai garis singgung tegak di $y_0$.
\end{theorem}

\begin{proof}
Untuk $y \to y_0$, tetapkanlah $x = f^{-1}(y)$: maka \omterm{def:b1:continuity:continuous}{kekontinuan} $f^{-1}$ memberikan $x
\to x_0$, dan
\[
\frac{f^{-1}(y) - f^{-1}(y_0)}{y - y_0}
= \frac{x - x_0}{f(x) - f(x_0)}
= \frac{1}{\dfrac{f(x) - f(x_0)}{x - x_0}}
\longrightarrow \frac{1}{f'(x_0)} .
\]
Adapun klaim garis singgung tegaknya: jika $f'(x_0) = 0$, maka hasil bagi yang terpampang itu
adalah kebalikan sebuah besaran yang menuju $0$ sambil mempertahankan
satu tanda yang tetap (karena untuk $f$ yang naik tegas, $\frac{f(x) -
f(x_0)}{x - x_0} > 0$ untuk setiap $x \neq x_0$): sehingga hasil bagi selisih
$f^{-1}$ menuju $+\infty$ (dan ke $-\infty$
untuk $f$ yang turun). Jadi inversnya tetap \omterm{def:b1:continuity:continuous}{kontinu} tetapi tak
\omterm{def:b1:derivative:def}{dapat diturunkan} di $y_0$ --- karena grafiknya, yaitu pencerminan grafik $f$
terhadap diagonalnya, berdiri tegak persis di tempat grafik $f$ berjalan
mendatar, sebagaimana diperlihatkan $x^{1/3}$ di $0$ terhadap $x^3$.
\end{proof}

\begin{example}[Turunan invers, dua kali]\label{ex:b1:derivative:inverseaction}
Teoremanya menghitung ulang \omterm{def:b1:derivative:def}{turunan} klasiknya tanpa kerja
limit apa pun. Untuk $\ln = \exp^{-1}$: di $y = \eu^x$,
\[
(\ln)'(y) = \frac{1}{\exp'(x)} = \frac{1}{\eu^{x}} = \frac1y ,
\]
yang sah untuk setiap $y > 0$ karena $\exp' = \exp$ tak pernah lenyap. Untuk
$\arctan = \tan^{-1}$: di $y = \tan x$,
\[
(\arctan)'(y) = \frac{1}{1 + \tan^2 x} = \frac{1}{1 + y^2} ,
\]
dengan memakai $\tan' = 1 + \tan^2 > 0$. Inti gagasan penutupnya: bahwa rumusnya
mengubah pengetahuan tentang sebuah fungsi menjadi pengetahuan tentang
inversnya dengan ongkos satu penyulihan --- dan
penyulihannya ($x = \ln y$, $x = \arctan y$) persis merupakan
\omterm{def:b1:logic:statement}{pernyataan} bahwa kedua variabelnya hidup di sisi yang berlawanan pada
bijeksinya.
\end{example}

\section{Rolle dan teorema nilai rata-rata}

\begin{proposition}[Ekstremum interior]\label{prop:b1:derivative:fermat}
Jika $f$ \omterm{def:b1:derivative:def}{dapat diturunkan} di sebuah titik \emph{\omterm{def:b1:topology:closure}{interior}} $x_0$ pada $I$ dan
mempunyai ekstremum lokal di sana, maka $f'(x_0) = 0$.
\end{proposition}

\begin{proof}
Katakanlah sebuah maksimum lokal: maka ada $r > 0$ dengan $f(x) \leq f(x_0)$
untuk $\abs{x - x_0} \leq r$, dan keinteriorannya menjamin bahwa kedua
sisi $x_0$ tersedia di dalam $I$. Untuk $0 < h \leq r$ maka
hasil bagi $\frac{f(x_0 + h) - f(x_0)}{h}$ berpembilang $\leq 0$
dan berpenyebut $> 0$: sehingga ia $\leq 0$, dan limitnya $f'(x_0)$
mewarisi $\leq 0$ (karena ketaksamaan yang lebar diteruskan ke limitnya,
\cref{thm:b1:seq:order}); sedangkan untuk $-r \leq h < 0$ hasil baginya
$\geq 0$, yang memberikan $f'(x_0) \geq 0$. Jadi $f'(x_0) = 0$. (Adapun di sebuah
titik ujung, hanya satu tandanya yang tersedia: sehingga kesimpulannya gagal di sana ---
pikirkanlah $x$ pada $\intcc{0}{1}$, yang maksimal di $1$ dengan \omterm{def:b1:derivative:def}{turunan}
$1$.)
\end{proof}

\begin{theorem}[Rolle]\label{thm:b1:derivative:rolle}
Misalkan $f$ \omterm{def:b1:continuity:continuous}{kontinu} pada $\intcc{a}{b}$, \omterm{def:b1:derivative:def}{dapat diturunkan} pada
$\intoo{a}{b}$, dengan $f(a) = f(b)$. Maka $f'(c) = 0$ untuk suatu $c \in
\intoo{a}{b}$.\index{teorema Rolle}
\end{theorem}

\begin{proof}
Menurut teorema nilai ekstrem (\cref{thm:b1:continuity:evt}), $f$
mencapai maksimum dan minimumnya pada $\intcc{a}{b}$. Jika keduanya
tercapai di titik ujungnya, maka (karena $f(a) = f(b)$) maks $=$ min sehingga $f$
konstan: jadi sebarang $c$ yang \omterm{def:b1:topology:closure}{interior} berlaku. Kalau tidak, sebuah ekstremumnya
tercapai di sebuah titik \omterm{def:b1:topology:closure}{interior} $c$, lalu
\cref{prop:b1:derivative:fermat} memberikan $f'(c) = 0$.
\end{proof}

\begin{theorem}[Teorema nilai rata-rata]\label{thm:b1:derivative:mvt}
Misalkan $f$ \omterm{def:b1:continuity:continuous}{kontinu} pada $\intcc{a}{b}$, \omterm{def:b1:derivative:def}{dapat diturunkan} pada
$\intoo{a}{b}$. Maka ada $c \in \intoo{a}{b}$ dengan
\[
f(b) - f(a) = f'(c)\,(b - a) .
\]
\emph{Ketaksamaan nilai rata-rata:} jika lebih lanjut $m \leq f' \leq M$ pada
$\intoo{a}{b}$, maka $m(b-a) \leq f(b) - f(a) \leq M(b-a)$;
khususnya $\abs{f'} \leq K$ mengakibatkan bahwa $f$ bersifat Lipschitz-$K$.
\index{teorema nilai rata-rata}
\end{theorem}

\begin{proof}
Terapkanlah Rolle pada $g(x) = f(x) - \frac{f(b) - f(a)}{b - a}(x - a)$: maka $g$
\omterm{def:b1:continuity:continuous}{kontinu} pada $\intcc{a}{b}$, \omterm{def:b1:derivative:def}{dapat diturunkan} di dalamnya, dan $g(a) =
f(a) = g(b)$. Di titik $c$ yang di situ $g'(c) = 0$: $f'(c) =
\frac{f(b)-f(a)}{b-a}$. Lalu ketaksamaannya menyusul dengan membatasi $f'(c)$;
sedangkan \omterm{def:b1:logic:statement}{pernyataan} Lipschitznya menerapkannya pada setiap pasang titiknya.
\end{proof}

\begin{omfigure}
\begin{tikzpicture}
\begin{axis}[omaxis, width=8.2cm, height=5.4cm,
    xmin=0, xmax=4.6, ymin=0, ymax=3.2,
    xtick={0.7,2.2,3.9}, xticklabels={$a$,$c$,$b$}, ytick=\empty]
  \addplot[omDef, very thick, domain=0.7:3.9, samples=100]
    {0.5 + 1.4*ln(x+0.3)};
  \addplot[omExo, thick] coordinates {(0.7,1.5) (3.9,2.5)};
  \addplot[omThm, thick, dashed, domain=1.1:3.4]
    {2.06 + 0.3125*(x - 2.2)};
  \fill[omThm] (axis cs:2.2,2.06) circle (1.6pt);
\end{axis}
\end{tikzpicture}

{\small Teorema nilai rata-rata: bahwa suatu garis singgung (yang putus-putus) sejajar dengan
tali busurnya (yang kelabu). Absisnya $c$ adalah tempat teorema Rolle, yang diterapkan
pada fungsinya dikurangi tali busurnya, menemukan sebuah titik kritis.}
\end{omfigure}

\begin{example}[Metode Newton itu metode Heron]\label{ex:b1:derivative:newton}
Metode Newton untuk memecahkan $f(x) = 0$ mengganti kurvanya dengan
garis singgungnya di tebakan yang sedang berlaku $x_n$ lalu mengambil akar garis
singgung itu sebagai tebakan berikutnya:
\[
0 = f(x_n) + f'(x_n)(x_{n+1} - x_n)
\quad\Longrightarrow\quad
x_{n+1} = x_n - \frac{f(x_n)}{f'(x_n)} .
\]
Jalankanlah ia pada $f(x) = x^2 - 2$:
\[
x_{n+1} = x_n - \frac{x_n^2 - 2}{2x_n}
= \frac{x_n}{2} + \frac{1}{x_n}
= \frac12\Bigl(x_n + \frac{2}{x_n}\Bigr) :
\]
yaitu persis iterasi Heron (\cref{ex:b1:seq:heron}), dua
milenium lebih awal. Adapun kecepatan kuadratik yang teramati di sana kini
terjelaskan oleh gambaran garis singgungnya: karena di dekat sebuah akar sederhana, kurva dan
garis singgungnya berselisih galat berorde dua, sehingga setiap langkahnya kira-kira
mengkuadratkan galatnya --- adapun \omterm{def:b1:logic:statement}{pernyataan} umumnya menyusul dari
batas Taylor pada \cref{ch:b1:taylor}. Inti gagasan penutupnya:
bahwa di tempat dikotomi (\cref{ex:b1:continuity:dichotomy}) hanya memakai
\omterm{def:b1:continuity:continuous}{kekontinuan} dan memperoleh satu bit per langkah, Newton membelanjakan sebuah
\omterm{def:b1:derivative:def}{turunan} untuk melipatduakan cacah angka yang benar per langkah.
\end{example}

\begin{example}[Teorema nilai rata-rata sebagai penaksir]\label{ex:b1:derivative:mvtestimate}
Seberapa besarkah $\sqrt{101}$? Terapkanlah teoremanya pada $f(t) = \sqrt t$
pada $\intcc{100}{101}$: maka untuk suatu $c \in \intoo{100}{101}$,
\[
\sqrt{101} - 10 = \frac{1}{2\sqrt c},
\qquad\text{sehingga}\qquad
\frac{1}{2\sqrt{101}} < \sqrt{101} - 10 < \frac{1}{20} = 0.05 ,
\]
dan karena $\sqrt{101} < 10.05$, batas kirinya melampaui
$\frac{1}{20.1} > 0.0497$: jadi $10.0497 < \sqrt{101} < 10.05$
(dengan nilai sejatinya $10.049875\dots$) --- yaitu tiga desimal yang benar dari satu
penilaian \omterm{def:b1:derivative:def}{turunan}. Demikian pula $\abs{\sin a - \sin b} \leq
\abs{a - b}$ (dengan batas $\abs{\cos}\leq 1$): jadi taksiran Lipschitz
yang dipakai sejak \cref{ch:b1:seq} semuanya teorema ini. Inti gagasan
penutupnya: bahwa teorema nilai rata-rata adalah rumus Taylor berorde nol
--- ia menukar satu titik $c$ yang tak diketahui dengan sebuah ketaksamaan yang tajam, dan
\cref{ch:b1:taylor} akan mengulangi pertukaran itu persis.
\end{example}

\begin{corollary}[Kriteria kemonotonan]\label{cor:b1:derivative:monotone}
Misalkan $f$ \omterm{def:b1:continuity:continuous}{kontinu} pada $I$, \omterm{def:b1:derivative:def}{dapat diturunkan} pada \omterm{def:b1:topology:closure}{interiornya}.
\begin{enumerate}
  \item $f' \geq 0$ pada \omterm{def:b1:topology:closure}{interiornya} $\iff$ $f$ naik; sedangkan $f' = 0$
        $\iff$ $f$ konstan.
  \item Jika $f' > 0$ kecuali di titik yang berhingga banyak yang di situ ia lenyap, maka
        $f$ naik \emph{tegas}.
\end{enumerate}
\end{corollary}

\begin{proof}
Jika $f' \geq 0$: maka untuk $x < y$ di $I$, teorema nilai rata-rata pada
$\intcc{x}{y}$ memberikan $f(y) - f(x) = f'(c)(y - x) \geq 0$. Sebaliknya,
hasil bagi selisih fungsi yang naik bernilai $\geq 0$, sehingga
limitnya pun demikian. Adapun kasus konstannya: terapkanlah yang sebelumnya pada $f'$ dan
$-f' \geq 0$. Untuk versi tegasnya: $f$ naik; dan kesamaan $f(x) =
f(y)$ untuk $x < y$ akan membekukan $f$ pada $\intcc{x}{y}$, yang memaksa $f' =
0$ di sana --- yaitu titik yang tak hingga banyaknya.
\end{proof}

\begin{example}[Turunan yang sama, fungsi yang berbeda]\label{ex:b1:derivative:intervalmatters}
Pada $\R^* = \intoo{-\infty}{0} \cup \intoo{0}{+\infty}$, baik
$f(x) = \ln\abs x$ maupun $g(x) = \ln\abs x + \mathbf{1}_{x>0}$
(dengan menambahkan $1$ hanya pada setengah garis kanannya) memenuhi $f' = g' =
\frac1x$. Namun keduanya tak berselisih konstanta: karena kriteria
``$f' = 0 \implies f$ konstan'' merupakan \omterm{def:b1:logic:statement}{pernyataan} tentang \emph{\omterm{prop:b1:reals:intervals}{selang}}
--- sebab buktinya menjalankan teorema nilai rata-rata antara dua titik,
yang menuntut seluruh ruas yang menghubungkannya terletak di
daerah asalnya. Pada masing-masing setengah garisnya secara tersendiri, primitif
$\frac1x$ adalah $\ln\abs x + c$ dengan satu konstanta per setengah garis,
jadi dua konstanta yang saling bebas seluruhnya. \Cref{ch:b1:integration}
mewarisi catatan halus ini: bahwa ``sang'' primitif sebuah fungsi
terdefinisi dengan baik sampai sebuah konstanta \emph{pada setiap \omterm{prop:b1:reals:intervals}{selang} daerah
asalnya}, dan tabel antiturunan diam-diam menganggap keterhubungannya.
\end{example}

\begin{example}[Ketegasan secara cuma-cuma]\label{ex:b1:derivative:strictness}
Fungsi $x \mapsto x^3$ naik \emph{tegas} pada $\R$ meskipun
\omterm{def:b1:derivative:def}{turunannya} lenyap di $0$: karena klausa kriterianya
``$f' > 0$ kecuali di titik yang berhingga banyak'' persis dirancang
untuk titik yang rata semacam itu. Sebaliknya, $f' \geq 0$ sendirian hanya memberikan
kenaikan dalam arti yang lebar (karena fungsi konstan pun memenuhinya), dan
\omterm{def:b1:derivative:def}{turunan} yang lenyap pada seluruh subselang memang membekukan
fungsinya di sana. Adapun aturan praktisnya: untuk mengklaim kemonotonan
yang tegas, daftarkanlah nol $f'$; berhingga banyak (atau lebih
umum, tak satu pun pada subselang mana pun) tak berbahaya, sedangkan sebuah \omterm{prop:b1:reals:intervals}{selang}
berisi nol bersifat mematikan.
\end{example}

\begin{example}[Telaah variasi yang lengkap]\label{ex:b1:derivative:variations}
Telaahlah $f(x) = x^3 - 3x + 1$ pada $\R$. \omterm{def:b1:derivative:def}{Turunannya}: $f'(x) =
3(x^2 - 1)$, yang positif pada $\intoo{-\infty}{-1}$, negatif pada
$\intoo{-1}{1}$, dan positif pada $\intoo{1}{+\infty}$: sehingga menurut
kriteria kemonotonannya, $f$ naik, lalu turun, lalu
naik, dengan maksimum lokal $f(-1) = 3$ dan minimum lokal
$f(1) = -1$. Limitnya: $\mp\infty$ di $\mp\infty$. Adapun akibatnya,
yang terbaca dari tabel variasinya dengan teorema nilai
antara pada setiap cabang monotonnya: bahwa $f$ lenyap tepat sekali pada
masing-masing
\[
\intoo{-\infty}{-1}, \qquad \intoo{-1}{1}, \qquad
\intoo{1}{+\infty}
\]
(karena nilai di sambungannya bertanda berlawanan: $3 > 0 > -1$),
sehingga persamaan $x^3 - 3x + 1 = 0$ mempunyai tepat tiga akar
real; dan secara numerik akarnya terletak di dekat $-1.88$, $0.35$, $1.53$. Inti
gagasan penutupnya: bahwa tabel variasi merupakan \emph{perkakas bukti},
bukan sketsa --- karena cabang monoton ditambah perubahan tanda sama dengan
tepat satu akar, dan tabelnya mendaftar cabangnya
secara menyeluruh.
\end{example}

\begin{theorem}[Rumus Leibniz]\label{thm:b1:derivative:leibniz}
Jika $f, g$ \omterm{def:b1:derivative:def}{dapat diturunkan} $n$ kali, maka $fg$ pun demikian,
dan\index{rumus Leibniz}
\[
(fg)^{(n)} = \sum_{k=0}^{n} \binom nk f^{(k)}\, g^{(n-k)} .
\]
\end{theorem}

\begin{proof}
Lewat induksi pada $n$, yang persis sejajar dengan teorema binomialnya. Adapun
kasus $n = 1$ merupakan aturan hasil kalinya. Dengan menganggap rumusnya berlaku pada peringkat
$n$, turunkanlah sekali lagi:
\[
(fg)^{(n+1)} = \sum_{k=0}^{n} \binom nk
\Bigl( f^{(k+1)} g^{(n-k)} + f^{(k)} g^{(n-k+1)} \Bigr),
\]
lalu indeks ulanglah jumlah pertamanya dengan $j = k + 1$ dan kumpulkanlah
koefisien $f^{(j)} g^{(n+1-j)}$: yaitu $\binom{n}{j-1} +
\binom nj = \binom{n+1}{j}$ menurut aturan Pascal
(\cref{prop:b1:counting:identities}), dengan suku \omterm{def:b1:topology:closure}{perbatasannya} $j = 0$
dan $j = n + 1$ membawa $\binom{n+1}{0} = \binom{n+1}{n+1} =
1$ sebagaimana seharusnya.
\end{proof}

\begin{example}[Leibniz dalam kerja]\label{ex:b1:derivative:leibnizaction}
Hitunglah $\bigl(x^2 \eu^x\bigr)^{(n)}$ untuk $n \geq 2$. Ambillah $f =
x^2$, yang \omterm{def:b1:derivative:def}{turunannya} mati cepat ($f' = 2x$, $f'' = 2$,
$f^{(k)} = 0$ untuk $k \geq 3$), dan $g = \eu^x$: maka hanya tiga suku
jumlah Leibniznya yang bertahan,
\[
\bigl(x^2\eu^x\bigr)^{(n)}
= \binom n0 x^2 \eu^x + \binom n1 (2x)\,\eu^x + \binom n2\,
2\,\eu^x
= \eu^x\bigl(x^2 + 2nx + n(n-1)\bigr).
\]
Periksa kewarasannya di $n = 1$: $\eu^x(x^2 + 2x)$, yang memang
$(x^2\eu^x)'$. Inti gagasan penutupnya: pakailah Leibniz ketika satu faktornya
sebuah \omterm{def:b1:poly:def}{polinomial} --- karena jumlahnya lalu hanya bersuku $\deg + 1$, dan
rumusnya menjadi bentuk \omterm{def:b1:topology:closed}{tertutup}, bukan kesamaan yang abstrak. (Adapun untuk dua
faktor yang sama-sama lincah tak hingga seperti $\eu^x\sin x$, eksponensial
kompleks dari \cref{ch:b1:complex} merupakan perkakas yang lebih baik.)
\end{example}

\section{Kecembungan}

\begin{definition}\label{def:b1:derivative:convex}
Fungsi $f \colon I \to \R$ disebut \emph{cembung}\index{fungsi cembung} bila
setiap tali busurnya terletak di atas grafiknya:
\[
\forall x, y \in I,\ \forall t \in \intcc{0}{1}, \quad
f\bigl(tx + (1-t)y\bigr) \leq t f(x) + (1-t) f(y).
\]
($f$ disebut \emph{cekung} bila $-f$ cembung.)
\end{definition}

\begin{theorem}[Pencirian diferensialnya]\label{thm:b1:derivative:convexchar}
Misalkan $f$ \omterm{def:b1:derivative:def}{dapat diturunkan} pada $I$. Berikut ini setara:
\begin{enumerate}
  \item $f$ \omterm{def:b1:derivative:convex}{cembung};
  \item $f'$ naik pada $I$;
  \item grafiknya terletak di atas setiap garis singgungnya: $f(y) \geq f(x) +
        f'(x)(y - x)$ untuk setiap $x, y \in I$.
\end{enumerate}
Jika $f$ \omterm{def:b1:derivative:def}{dapat diturunkan} dua kali: maka $f$ \omterm{def:b1:derivative:convex}{cembung} $\iff f'' \geq 0$.
\end{theorem}

\begin{proof}
(1 $\Rightarrow$ 3) Kecembungannya yang ditulis sebagai $\frac{f(x + t(y-x)) -
f(x)}{t} \leq f(y) - f(x)$ untuk $t \in \intoc{0}{1}$; lalu biarkan $t \to 0^+$:
$f'(x)(y - x) \leq f(y) - f(x)$.

(3 $\Rightarrow$ 2) Untuk $x < y$, kedua ketaksamaan garis singgungnya di $x$
dan di $y$ memberikan $f'(x)(y-x) \leq f(y) - f(x) \leq f'(y)(y - x)$,
sehingga $f'(x) \leq f'(y)$.

(2 $\Rightarrow$ 1) Tetapkanlah $x < y$ dan $t \in \intoo{0}{1}$, lalu misalkan $z
= tx + (1-t)y \in \intoo{x}{y}$. Menurut teorema nilai rata-rata pada
$\intcc{x}{z}$ dan $\intcc{z}{y}$: ada $c_1 < z < c_2$ dengan
\[
\frac{f(z) - f(x)}{z - x} = f'(c_1) \leq f'(c_2)
= \frac{f(y) - f(z)}{y - z} ,
\]
lalu membersihkan penyebutnya ($z - x = (1-t)(y-x)$, $y - z = t(y-x)$)
menyusunnya ulang persis menjadi ketaksamaan kecembungannya.

Adapun kasus \omterm{def:b1:derivative:def}{dapat diturunkan} dua kalinya: $f'' \geq 0 \iff f'$ naik
(\cref{cor:b1:derivative:monotone}).
\end{proof}

\begin{omfigure}
\begin{tikzpicture}
\begin{axis}[omaxis, width=8.2cm, height=5.6cm,
    xmin=-2.4, xmax=2.6, ymin=-0.8, ymax=3.4,
    xtick=\empty, ytick=\empty]
  \addplot[omDef, very thick, domain=-2.1:2.3, samples=100]
    {0.5*x*x};
  \addplot[omExo, thick] coordinates {(-1.6,1.28) (2,2)};
  \addplot[omThm, thick, dashed, domain=-1.2:2.4] {0.72*(x-0.72)+0.2592};
\end{axis}
\end{tikzpicture}

{\small Kecembungan, dua kali: bahwa setiap tali busurnya (yang kelabu) terletak di atas grafiknya,
dan grafiknya terletak di atas setiap garis singgungnya (yang putus-putus).}
\end{omfigure}

\begin{example}[Ketaksamaan kecembungan yang klasik]\label{ex:b1:derivative:convexineq}
Fungsi $\exp$ \omterm{def:b1:derivative:convex}{cembung} (karena $\exp'' = \exp > 0$): sehingga garis singgungnya di $0$ memberikan
$\eu^x \geq 1 + x$ untuk setiap $x$. Sedangkan $\ln$ cekung: sehingga garis singgungnya di $1$
memberikan $\ln x \leq x - 1$; dan tali busurnya memberikan, untuk $0 < a \leq b$,
ketaksamaan antara rata-rata geometri dan aritmetika: dengan mengambil $t =
\frac12$ pada kecekungannya,
\[
\ln\frac{a + b}{2} \geq \frac{\ln a + \ln b}{2} = \ln\sqrt{ab},
\qquad\text{sehingga}\qquad
\sqrt{ab} \leq \frac{a+b}{2} .
\]
Adapun ketaksamaan aritmetika--geometri yang umum adalah
\cref{exo:b1:derivative:9}.
\end{example}

\begin{example}[Sebuah ketaksamaan kecembungan dari nol]\label{ex:b1:derivative:xlnx}
Fungsi $f(t) = t\ln t$ bersifat \omterm{def:b1:derivative:convex}{cembung} pada $\intoo{0}{+\infty}$:
karena $f''(t) = \frac1t > 0$. Adapun ketaksamaan titik tengahnya, yang dikalikan
$2$, berbunyi: untuk setiap $a, b > 0$,
\[
a\ln a + b\ln b \;\geq\; (a + b)\,\ln\frac{a + b}{2} ,
\]
dengan kesamaannya jika dan hanya jika $a = b$ (menurut kecembungan tegasnya). Uji cobanya: $a =
1$, $b = 3$ memberikan $3\ln 3 = 3.296$ terhadap $4\ln 2 = 2.773$.
Ketaksamaan yang tampak polos ini merupakan kasus dua titik pada
perbandingan \emph{entropi} yang muncul kembali bersama ketaksamaan
Jensen (\cref{exo:b1:derivative:9}) dan pada
asimtotik teori informasi jilid Tahun ke-3. Inti
gagasan penutupnya: bahwa untuk memproduksi sebuah ketaksamaan, carilah fungsi
yang \omterm{def:b1:derivative:def}{turunan} keduanya bertanda lalu tuliskanlah apa yang dikatakan kecembungannya
--- karena pencirian diferensialnya mengubah satu pemeriksaan tanda
menjadi ketaksamaan yang tak hingga banyaknya.
\end{example}

\begin{remark}[Jebakan yang lazim dengan turunan]
(i) \emph{\omterm{def:b1:derivative:def}{Turunan} yang positif di satu titik tak memberikan
kemonotonan di dekatnya}: karena $f(x) = \frac x2 + x^2\sin\frac1x$ (dengan
$f(0) = 0$) mempunyai $f'(0) = \frac12 > 0$, namun
\[
f'(x) = \frac12 + 2x\sin\frac1x - \cos\frac1x
\]
sama dengan $-\frac12$ di setiap $x_n = \frac{1}{2\pi n}$: jadi setiap
\omterm{def:b1:topology:open}{persekitaran} $0$ memuat penurunan. Adapun kemonotonan menuntut $f'
\geq 0$ \emph{pada sebuah \omterm{prop:b1:reals:intervals}{selang}}
(\cref{cor:b1:derivative:monotone}); sedangkan tanda titik demi titiknya hanya
mengendalikan perpotongan dengan garis singgungnya. (ii) \emph{Ketiga
hipotesis Rolle semuanya aktif}: karena $\abs x$ pada $\intcc{-1}{1}$
(tanpa sifat \omterm{def:b1:derivative:def}{dapat diturunkan} \omterm{def:b1:topology:closure}{interior}), $x$ pada $\intcc{0}{1}$ (dengan ujung yang tak
sama), dan $x - \lfloor x\rfloor$ pada $\intcc{0}{1}$ (yang \omterm{def:b1:continuity:continuous}{kekontinuannya}
gagal di $1$) masing-masing mematahkan tepat satu hipotesisnya beserta
kesimpulannya. (iii) \emph{\omterm{def:b1:derivative:def}{Turunan} boleh takkontinu, tetapi
tak sebarangan}: karena $f'$ dapat berayun
(\cref{ex:b1:derivative:notc1}) namun selalu memenuhi
sifat nilai antara (menurut Darboux,
\cref{exo:b1:derivative:10}): jadi sebuah \omterm{def:b1:derivative:def}{turunan} tak pernah melompat --- sehingga jika
kamu menghitung sebuah ``limit \omterm{def:b1:derivative:def}{turunan}'' sepihak yang melompat, maka kamu
sudah menurunkan fungsi yang tak \omterm{def:b1:derivative:def}{dapat diturunkan}. (iv) \emph{Rumus
inversnya menuntut $f' \neq 0$}: karena $x \mapsto x^3$ merupakan bijeksi mulus
yang naik tegas yang inversnya $x^{1/3}$ mempunyai
garis singgung tegak di $0$ --- jadi sifat \omterm{def:b1:derivative:def}{dapat diturunkan} inversnya
hilang persis di tempat $f'$ lenyap
(\cref{thm:b1:derivative:inverse}).
\end{remark}

\begin{remark}[Di mana teorema nilai rata-rata bekerja berikutnya]
Hampir setiap \omterm{def:b1:logic:statement}{pernyataan} kuantitatif pada bab berikutnya merupakan
teorema nilai rata-rata bab ini yang berkostum: teorema fundamental
kalkulus (\cref{ch:b1:integration}) menurunkan fungsi
luasnya lalu menutup dengan kriteria kemonotonannya; rumus
Taylor--Lagrange (\cref{ch:b1:taylor}) merupakan teorema nilai rata-rata
yang diulang $n$ kali; analisis galat metode Newton
dan iterasi titik tetap (\cref{exo:b1:derivative:11}) merupakan
bentuk Lipschitznya; sedangkan soal akhir pekan bab ini
(\cref{pb:b1:derivative:1}) mengubah batas Lipschitz yang sama menjadi
teori bilangan --- yaitu ketaksamaan tolakan antara \omterm{pb:b1:logic:1}{bilangan aljabar}
dan bilangan rasional, yang menghasilkan \omterm{pb:b1:logic:1}{bilangan transenden} yang pertama
dalam sejarah. Pada jilid Tahun ke-2, \emph{ketaksamaan} nilai
rata-ratanya bertahan dalam beberapa variabel ketika kesamaannya
tidak.
\end{remark}

\begin{example}[Ketaksamaan Young dari kecekungan]\label{ex:b1:derivative:young}
Misalkan $p, q > 1$ dengan $\frac1p + \frac1q = 1$. Untuk setiap $a, b > 0$:
\[
ab \;\leq\; \frac{a^p}{p} + \frac{b^q}{q} .
\]
Buktinya lewat satu penerapan kecekungan $\ln$ dengan bobot
$\frac1p, \frac1q$ (yaitu ketaksamaan Jensen dua titik, seperti pada
\cref{exo:b1:derivative:9}):
\[
\ln\Bigl(\frac{a^p}{p} + \frac{b^q}{q}\Bigr)
\;\geq\; \frac1p \ln(a^p) + \frac1q \ln(b^q)
= \ln a + \ln b = \ln(ab),
\]
lalu $\ln$ yang naik mengubah ketaksamaan logaritmanya menjadi
klaimnya; dengan kesamaannya jika dan hanya jika $a^p = b^q$ (menurut kecekungan tegasnya). Adapun kasus
$p = q = 2$ merupakan ketaksamaan aritmetika-geometri $ab \leq
\frac{a^2 + b^2}{2}$ yang menyamar. Inti gagasan penutupnya: bahwa ketaksamaan
Young merupakan benih aljabar bagi ketaksamaan Hölder dan Minkowski
pada jilid Tahun ke-2 --- yaitu satu \omterm{def:b1:logic:statement}{pernyataan} kecekungan
tentang $\ln$, yang dipanen bagi norma.
\end{example}

\begin{remark}[Cakrawala di dalam jilid ini]
\omterm{def:b1:derivative:def}{Turunan} memperoleh tiga kehidupan baru sebelum jilid ini berakhir.
Pada \cref{ch:b1:taylor} ia mengulang dirinya: $n$ \omterm{def:b1:derivative:def}{turunan} di sebuah titik
memampat menjadi satu \omterm{def:b1:poly:def}{polinomial} ditambah galat yang terkendali, dan
teorema nilai rata-ratanya menjadi sisa Lagrange. Pada
\cref{ch:b1:curves}, pendiferensialan menjadi geometris: karena untuk
kurva berparameter $t \mapsto (x(t), y(t))$, pasangan $(x'(t),
y'(t))$ itu adalah \emph{vektor} kecepatan, ketersinggungannya menjadi
kesegarisan, dan titik kritisnya menjadi katup yang harus digolongkan. Pada
\cref{ch:b1:multivar}, satu variabelnya dibekukan pada satu waktu: sehingga \omterm{def:b1:derivative:def}{turunan}
parsialnya mengulangi bab ini dua kali lipat, dan garis singgungnya
tumbuh menjadi bidang singgung. Ketiga babnya mewarisi tata bahasa
yang sama --- hampiran linear yang lokal ditambah suku galatnya ---
yang mula-mula diucapkan di sini.
\end{remark}

\section{Latihan}

\begin{exercise}[$\star$]\label{exo:b1:derivative:1}
Turunkanlah (dengan menyebutkan daerah asalnya): $x^x$;
$\;\ln\bigl(x + \sqrt{x^2+1}\bigr)$; $\;\arctan\frac{1}{x}$;
$\;\sqrt{1 + \eu^{2x}}$.
\end{exercise}

\begin{exercise}[$\star$]\label{exo:b1:derivative:2}
Lengkapilah \cref{ex:b1:derivative:notc1}: buktikan bahwa $f(x) = x^2
\sin\frac1x$, $f(0) = 0$, \omterm{def:b1:derivative:def}{dapat diturunkan} di $0$ dengan $f'(0) = 0$,
dan bahwa $f'$ tak berlimit di $0$.
\end{exercise}

\begin{exercise}[$\star$]\label{exo:b1:derivative:3}
Dengan memakai teorema nilai rata-rata atau ketaksamaan garis singgungnya, buktikan bahwa
untuk setiap $x > 0$:
\[
\frac{x}{1 + x} < \ln(1 + x) < x .
\]
Lalu simpulkan $\lim_{n\to\infty} \bigl(1 + \frac xn\bigr)^n = \eu^x$ untuk
setiap $x > 0$.
\end{exercise}

\begin{exercise}[$\star$]\label{exo:b1:derivative:4}
Misalkan $P$ \omterm{def:b1:poly:def}{polinomial} real dengan $k$ akar real yang berbeda. Buktikan bahwa
$P'$ mempunyai sekurang-kurangnya $k - 1$ akar real yang berbeda, yang terselang-seling dengan akar
$P$. Lalu simpulkan bahwa jika semua akar $P$ real, maka akar $P'$ pun demikian.
\end{exercise}

\begin{exercise}[$\star\star$]\label{exo:b1:derivative:5}
Misalkan $f$ \omterm{def:b1:derivative:def}{dapat diturunkan} pada $\R$ dengan $f' (x)\to \ell$ ketika $x \to
+\infty$. Buktikan bahwa $\frac{f(x)}{x} \to \ell$ \emph{(dengan teorema nilai
rata-rata pada $\intcc{A}{x}$)}. Apakah $f(x+1) - f(x) \to \ell$ juga berlaku?
\end{exercise}

\begin{exercise}[$\star\star$]\label{exo:b1:derivative:6}
(Rolle yang diskret) Misalkan $f$ \omterm{def:b1:derivative:def}{dapat diturunkan} $n$ kali pada $I$ dan
lenyap di $n + 1$ titik yang berbeda. Buktikan bahwa $f^{(n)}$ lenyap
sekurang-kurangnya sekali. Penerapannya: bahwa \omterm{def:b1:poly:def}{polinomial} berderajat $\leq n$ yang lenyap di
$n+1$ titik bernilai nol (sekali lagi).
\end{exercise}

\begin{exercise}[$\star\star$]\label{exo:b1:derivative:7}
Misalkan $f$ \omterm{def:b1:derivative:def}{dapat diturunkan} dua kali pada $\intcc{a}{b}$ dengan $f(a) = f(b)
= 0$ dan $f(x_0) > 0$ untuk suatu $x_0$ yang \omterm{def:b1:topology:closure}{interior}. Buktikan bahwa $f''(c) <
0$ untuk suatu $c \in \intoo{a}{b}$. \emph{(Dua teorema nilai rata-rata dan
satu perbandingan kemiringan.)}
\end{exercise}

\begin{exercise}[$\star\star$]\label{exo:b1:derivative:8}
Telaahlah fungsi $f(x) = \dfrac{\ln x}{x}$ pada
$\intoo{0}{+\infty}$: yaitu variasi, limit, dan maksimumnya. Simpulkan bahwa $a^b
> b^a$ untuk setiap bilangan real $\eu \leq a < b$, lalu selesaikanlah kasus khusus
yang termasyhur: manakah yang lebih besar di antara $\eu^\pi$ dan $\pi^\eu$? Periksalah terhadap
pasangan bilangan bulat kecil $(2,3)$ dan $(2,4)$: mengapa keduanya berperilaku
berbeda?
\end{exercise}

\begin{exercise}[$\star\star$]\label{exo:b1:derivative:9}
(Ketaksamaan aritmetika--geometri) Dengan memakai kecekungan $\ln$ dengan
bobot yang umum (yaitu ketaksamaan Jensen bagi $n$ titik, yang harus dibuktikan lewat
induksi pada $n$), tunjukkan bahwa untuk bilangan real positif $a_1, \dots, a_n$:
\[
\sqrt[n]{a_1 a_2 \cdots a_n} \leq \frac{a_1 + \dots + a_n}{n},
\]
dengan kesamaannya jika dan hanya jika semua $a_i$ sama.
\end{exercise}

\begin{exercise}[$\star\star\star$]\label{exo:b1:derivative:10}
(Darboux: \omterm{def:b1:derivative:def}{turunan} mengambil nilai antaranya) Misalkan $f$
\omterm{def:b1:derivative:def}{dapat diturunkan} pada $I$ dan $a < b$ di $I$ dengan $f'(a) < v < f'(b)$.
Dengan meninjau $g(x) = f(x) - vx$ dan titik yang di situ $g$ mencapai
minimumnya pada $\intcc{a}{b}$, buktikan bahwa $f'(c) = v$ untuk suatu $c
\in \intoo{a}{b}$ --- meskipun $f'$ tak harus \omterm{def:b1:continuity:continuous}{kontinu}
(\cref{exo:b1:derivative:2}).
\end{exercise}

\begin{exercise}[$\star\star\star$]\label{exo:b1:derivative:11}
Misalkan $f \colon \R \to \R$ \omterm{def:b1:derivative:def}{dapat diturunkan} dengan $\abs{f'(x)} \leq k
< 1$ untuk setiap $x$ (yaitu sebuah \emph{kontraksi}). Buktikan bahwa $f$ mempunyai tepat
satu titik tetap $\ell$, dan bahwa setiap barisan $u_{n+1} = f(u_n)$
konvergen ke $\ell$ dengan $\abs{u_n - \ell} \leq k^n\abs{u_0 -
\ell}$. \emph{(Untuk keberadaannya: terapkan teorema nilai antara pada
$f(x) - x$ pada sebuah ruas yang besar, dengan memakai batas Lipschitznya; atau pakailah
kelengkapannya dengan kriteria Cauchy.)}
\end{exercise}

\begin{exercise}[$\star\star\star$]\label{exo:b1:derivative:12}
(Teorema nilai rata-rata Cauchy dan aturan l'Hospital)
\begin{enumerate}
  \item Misalkan $f, g$ \omterm{def:b1:continuity:continuous}{kontinu} pada $\intcc{a}{b}$, dan
        \omterm{def:b1:derivative:def}{dapat diturunkan} pada $\intoo{a}{b}$, dengan $g'$ tak pernah nol
        di sana. Buktikan bahwa $g(b) \neq g(a)$ dan bahwa suatu $c \in
        \intoo{a}{b}$ memenuhi
        \[
        \frac{f(b) - f(a)}{g(b) - g(a)} = \frac{f'(c)}{g'(c)}
        \]
        \emph{(terapkanlah Rolle pada $h = f - \lambda g$ untuk konstanta
        $\lambda$ yang tepat)}.
  \item Simpulkanlah aturan l'Hospital dalam bentuk $\frac00$ di sebuah
        titik: bahwa jika $f(a) = g(a) = 0$ dan
        $\frac{f'(x)}{g'(x)} \to \ell$ ketika $x \to a^+$, maka
        $\frac{f(x)}{g(x)} \to \ell$.
  \item Tunjukkan bahwa konversnya gagal: bahwa untuk $f(x) = x^2\sin\frac1x$
        ($f(0) = 0$) dan $g(x) = x$, hasil bagi
        $\frac{f}{g}$ mempunyai limit di $0$ tetapi
        $\frac{f'}{g'}$ tak mempunyainya.
\end{enumerate}
\end{exercise}

\section{Soal: Ketaksamaan Liouville dan bilangan transenden yang
pertama}

\begin{problem}[{Soal akhir pekan --- \omterm{pb:b1:logic:1}{bilangan aljabar} menolak
bilangan rasional: $\abs{x - p/q} \geq C/q^d$, dan ketransendenan
$\sum 10^{-n!}$}]
\label{pb:b1:derivative:1}
Sebuah bilangan real disebut \emph{aljabar} bila ia akar sebuah
\omterm{def:b1:poly:def}{polinomial} taknol berkoefisien bilangan bulat, dan \emph{transenden}
bila tidak. Pada tahun 1844 Liouville menghasilkan bilangan pertama yang pernah
\emph{terbukti} transenden, dan mesin buktinya adalah
teorema nilai rata-rata bab ini: bahwa \omterm{pb:b1:logic:1}{bilangan aljabar} berderajat $d$
tak dapat dihampiri oleh bilangan rasional lebih baik daripada $C/q^d$ --- sehingga
bilangan yang terhampiri \emph{lebih cepat daripada setiap pangkat} tak dapat
bersifat aljabar. Soal ini membangun ketaksamaannya, mengonstruksi
bilangan Liouville $L = 0.110001000\dots$ (dengan angka satu pada posisi
faktorialnya, lewat mesin angka pada \cref{pb:b1:reals:1}),
membuktikan ketransendenannya, lalu berakhir dengan bukti tandingan Cantor dan
batas efektif bagi $\sqrt2$ dan
$2^{1/3}$.\index{bilangan transenden}\index{ketaksamaan
Liouville}\index{bilangan aljabar}

\medskip
\noindent\textbf{Bagian I --- Seberapa baik bilangan rasional dapat
dihampiri?}
\begin{enumerate}
  \item Tunjukkan bahwa dua bilangan rasional yang berbeda $\frac ab \neq \frac pq$
        (yang ditulis dengan $b, q \geq 1$) memenuhi $\bigl|\frac ab -
        \frac pq\bigr| \geq \frac{1}{bq}$. Lalu simpulkan: bahwa jika $x =
        \frac ab$ dan $0 < \bigl|x - \frac pq\bigr| <
        \frac{1}{bq}$, maka $\frac pq$ yang demikian tak ada --- jadi bilangan rasional
        menolak semua bilangan rasional lain pada skala $\frac 1q$.
  \item Buktikan bahwa untuk \emph{setiap} bilangan rasional $\frac pq$ (dengan $q \geq
        1$): $\bigl|\sqrt2 - \frac pq\bigr| \geq
        \frac{1}{4q^2}$ \emph{(jika jaraknya melampaui $1$ maka ini
        jelas; kalau tidak, batasilah $\abs{\sqrt2 + p/q} < 4$ lalu
        pakailah bilangan bulat taknol $\abs{p^2 - 2q^2} \geq 1$)}.
  \item Pada arah yang lain: periksalah bahwa $(p, q) \mapsto (p +
        2q, p + q)$ mengawetkan $\abs{p^2 - 2q^2} = 1$, lalu bangkitkanlah
        dari $(1,1)$ pasangan $(3,2)$, $(7,5)$, $(17,12)$,
        $(41,29)$, $(99,70)$, dan tunjukkan bahwa masing-masingnya memenuhi
        \[
        \Bigl|\sqrt2 - \frac pq\Bigr| =
        \frac{1}{q^2\,(\sqrt2 + p/q)} < \frac{1}{2q^2} :
        \]
        jadi hampiran berorde $2$ yang tak hingga banyaknya. Digabung dengan
        pertanyaan 2: eksponen hampiran $\sqrt 2$ \emph{tepat}
        sama dengan $2$.
  \item (Dirichlet) Misalkan $x$ irasional dan $N \in \N^*$.
        Tinjaulah $N + 1$ bagian pecahan $0, x, 2x,
        \dots, Nx$ pada $N$ kotak $\intco{\frac kN}{\frac{k +
        1}{N}}$: maka menurut prinsip sarang merpati
        (\cref{cor:b1:counting:pigeonhole}), dua di antaranya jatuh pada satu kotak.
        Simpulkanlah adanya $q \leq N$ dan $p$ dengan $\abs{qx - p} <
        \frac 1N$, sehingga ada bilangan rasional yang tak hingga banyaknya dengan
        $\bigl|x - \frac pq\bigr| < \frac{1}{q^2}$: jadi \emph{setiap}
        bilangan irasional terhampiri sampai orde $2$.
\end{enumerate}

\medskip
\noindent\textbf{Bagian II --- Ketaksamaan Liouville.} Misalkan $x$
irasional dan aljabar.
\begin{enumerate}[resume]
  \item Tunjukkan bahwa di antara \omterm{def:b1:poly:def}{polinomial} bulat yang taknol
        yang lenyap di $x$ ada satu, katakanlah $P$ berderajat $d$,
        yang \emph{tak berakar rasional}; lalu periksalah $d \geq 2$
        \emph{(bagilah keluar sebuah faktor $X - \frac ab$ atas $\Q$ lalu
        bersihkanlah penyebutnya; adapun derajat $1$ akan membuat $x$
        rasional)}.
  \item Tunjukkan bahwa untuk setiap bilangan rasional $\frac pq$ (dengan $q \geq 1$):
        $\bigl|P\bigl(\frac pq\bigr)\bigr| \geq \frac{1}{q^d}$
        \emph{(karena $q^d P(p/q)$ merupakan bilangan bulat taknol)}.
  \item Misalkan $M = \max_{\intcc{x-1}{x+1}} \abs{P'}$
        (\cref{thm:b1:continuity:evt}). Dengan memakai teorema nilai
        rata-rata antara $x$ dan $\frac pq$, buktikanlah
        \emph{ketaksamaan Liouville}: bahwa dengan $C = \min\bigl(1,
        \frac 1M\bigr) > 0$,
        \[
        \Bigl| x - \frac pq \Bigr| \geq \frac{C}{q^{\,d}}
        \qquad\text{untuk setiap bilangan rasional } \frac pq,\ q \geq 1 .
        \]
  \item Sebutlah $x$ sebuah \emph{bilangan Liouville} bila untuk setiap $n \in
        \N$ ada bilangan rasional $\frac pq$ dengan $q \geq 2$ dan
        $0 < \bigl|x - \frac pq\bigr| < q^{-n}$. Buktikan bahwa
        bilangan Liouville bersifat irasional \emph{(lewat pertanyaan 1: pilihlah
        $n$ dengan $2^{\,n-1} > b$)}.
  \item Buktikanlah teorema Liouville: bahwa \emph{bilangan Liouville bersifat
        transenden} \emph{(gabungkanlah pertanyaan 7 dan 8: karena
        ketaksamaan $C < q^{\,d-n}$ gagal untuk $n$ yang besar)}.
\end{enumerate}

\medskip
\noindent\textbf{Bagian III --- Bilangan $L$.}
\begin{enumerate}[resume]
  \item Misalkan $L$ nilai (dalam pengertian
        \cref{pb:b1:reals:1}) untaian angka desimal yang berangka
        $1$ pada posisi $n!$ ($n = 1, 2, 3, \dots$)
        dan $0$ di tempat lain, yakni $L = \sup_k t_k$ dengan $t_k =
        \sum_{n=1}^{k} 10^{-n!}$. Tuliskanlah $25$ angka
        pertama $L$.
  \item Buktikan apitan ekornya, untuk setiap $k \geq 1$:
        \[
        10^{-(k+1)!} \;\leq\; L - t_k \;\leq\;
        \frac{10}{9}\,10^{-(k+1)!} \;<\; 2\cdot 10^{-(k+1)!}
        \]
        \emph{(batasilah setiap jumlah parsial di luar $t_k$ oleh jumlah
        geometri yang hingga)}.
  \item Tulislah $t_k = \frac{p_k}{q_k}$ dengan $q_k = 10^{k!}$.
        Tunjukkan $0 < L - \frac{p_k}{q_k} < \frac{2}{q_k^{\,k+1}}$,
        lalu simpulkan bahwa $L$ merupakan bilangan Liouville dalam pengertian
        pertanyaan 8.
  \item Simpulkan: bahwa $L$ transenden --- yaitu contoh eksplisit
        yang pertama dalam sejarah (Liouville, 1844). Periksalah silang
        keirasionalannya secara langsung: bahwa angkanya tak periodik pada
        akhirnya (karena jaraknya tumbuh, seperti pada \cref{pb:b1:reals:1},
        pertanyaan 20).
  \item Samaratakanlah: gantilah setiap angka $1$ dengan sebarang
        angka taknol $d_n \in \intint{1}{9}$. Tunjukkan bahwa nilainya
        tetap bilangan Liouville, lalu simpulkan --- lewat argumen diagonal
        pada \cref{pb:b1:reals:1} (pertanyaan 22) yang diterapkan
        pada pemilihan angka ini --- bahwa ada \omterm{pb:b1:logic:1}{bilangan transenden}
        berbentuk demikian yang takterbilang banyaknya.
\end{enumerate}

\medskip
\noindent\textbf{Bagian IV --- Hierarki orde
hampirannya.} Katakanlah $x$ \emph{terhampiri sampai orde $\mu$} bila untuk
suatu konstanta $c > 0$ ada bilangan rasional yang tak hingga banyaknya yang memenuhi
$\bigl|x - \frac pq\bigr| < \frac{c}{q^{\mu}}$.
\begin{enumerate}[resume]
  \item Rakitlah hierarkinya dari Bagian I--III: bahwa bilangan rasional
        terhampiri sampai orde $1$ dan tak lebih baik; $\sqrt 2$ sampai
        orde $2$ dan tak lebih baik; setiap bilangan irasional sampai orde
        sekurang-kurangnya $2$; \omterm{pb:b1:logic:1}{bilangan aljabar} berderajat $d$ tak sampai orde
        mana pun di luar $d$; sedangkan bilangan Liouville sampai setiap orde. Benarkanlah
        setiap klaimnya dengan mengutip pertanyaan yang relevan.
  \item Tunjukkan bahwa $L + r$ merupakan bilangan Liouville untuk setiap
        bilangan rasional $r = \frac ab$ \emph{(translasikanlah
        penghampirnya: karena penyebut barunya menjadi $b\,q_k$)}.
        Lalu simpulkan bahwa bilangan Liouville --- sehingga \omterm{pb:b1:logic:1}{bilangan transenden} ---
        bersifat \omterm{def:b1:topology:dense}{padat} di $\R$.
  \item (Cantor, 1874) Buktikan bahwa \omterm{def:b1:logic:sets}{himpunan} \omterm{pb:b1:logic:1}{bilangan aljabar}
        bersifat terbilang: karena ada berhingga banyak \omterm{def:b1:poly:def}{polinomial}
        bulat yang derajatnya ditambah jumlah $\abs{\text{koefisien}}$-nya
        terbatas oleh $h$, dan masing-masingnya berakar paling banyak $\deg$; sedangkan
        gabungan terbilang atas \omterm{def:b1:counting:card}{himpunan hingga} bersifat terbilang. Dan karena tak ada
        barisan yang menghabiskan $\R$ (\cref{pb:b1:reals:1}, pertanyaan
        22), maka \omterm{pb:b1:logic:1}{bilangan transenden} ada --- bahkan membentuk \omterm{def:b1:logic:sets}{himpunan}
        yang takterbilang. Bandingkanlah kedua buktinya: apakah yang diberikan
        bukti Liouville yang tak dapat diberikan bukti Cantor?
  \item Buktikan langsung dari pertanyaan 2 bahwa $\sqrt 2$
        \emph{bukan} bilangan Liouville \emph{(karena untuk $n \geq 3$,
        ketaksamaan $q^{-n} > \frac{1}{4q^2}$ membatasi $q$;
        lalu hanya berhingga banyak calon bilangan rasional yang tersisa, yang semuanya
        berjarak positif dari $\sqrt2$)}. Samaratakanlah:
        bahwa tak ada \omterm{pb:b1:logic:1}{bilangan aljabar} yang Liouville.
\end{enumerate}

\medskip
\noindent\textbf{Bagian V --- Konstanta yang efektif.}
\begin{enumerate}[resume]
  \item Untuk pasangan Pell $(99, 70)$: periksalah $99^2 - 2\cdot70^2
        = 1$ lalu nilailah galatnya yang persis
        \[
        \sqrt2 - \frac{99}{70}
        = \frac{-1}{70^2\,\bigl(\sqrt2 + \frac{99}{70}\bigr)},
        \qquad
        \Bigl|\sqrt 2 - \frac{99}{70}\Bigr| \approx 7.2\cdot
        10^{-5} :
        \]
        jadi lima angka yang benar dari pecahan berangka tiga.
  \item Jalankanlah Bagian II pada $x = 2^{1/3}$, $P = X^3 - 2$: periksalah bahwa $P$
        tak berakar rasional, lalu batasilah $M = \max_{\intcc{x-1}{x+1}}
        3t^2 \leq 3\,(1 + 2^{1/3})^2 < 16$, lalu simpulkan
        ketaksamaan efektifnya
        \[
        \Bigl| 2^{1/3} - \frac pq \Bigr| \geq
        \frac{1}{16\,q^3} \qquad \text{untuk setiap } \frac pq .
        \]
  \item Hasilnya: tunjukkan bahwa sebarang bilangan rasional yang menghampiri $2^{1/3}$
        dalam jarak $10^{-6}$ pasti berpenyebut $q \geq 40$.
  \item Tunjukkan bahwa basis $10$ tak relevan: bahwa padanan
        binernya $\sum_{n\geq1} 2^{-n!}$ (yaitu nilai untaian
        biner dengan angka satu pada posisi faktorialnya) juga merupakan
        bilangan Liouville, sehingga transenden.
\end{enumerate}

\medskip
\noindent\textbf{Bagian VI --- \omterm{def:b1:topology:closure}{Perbatasan} dan sintesisnya.}
\begin{enumerate}[resume]
  \item Misalkan $x^\dagger$ nilai untaian desimal dengan
        angka satu tepat pada posisi $3^k$ ($k \geq 0$). Tunjukkan bahwa
        $x^\dagger$ terhampiri sampai orde $3$, lalu simpulkan
        dari ketaksamaan Liouville bahwa $x^\dagger$ bukan
        bilangan rasional maupun bilangan irasional kuadratik. Jelaskanlah mengapa
        metodenya mandek di situ: karena orde $3$ selaras dengan
        kealjabaran berderajat $\geq 3$, dan menutup celah itu
        (karena sebarang eksponen $> 2$ sudah mencukupi, bagi setiap \omterm{pb:b1:logic:1}{bilangan
        aljabar}) merupakan teorema Roth, yang jauh di luar jilid ini.
  \item Kuantifikasikanlah Cantor: tunjukkan bahwa \omterm{pb:b1:logic:1}{bilangan aljabar}
        berderajat $\leq d$ yang diberikan oleh \omterm{def:b1:poly:def}{polinomial} berkoefisien
        di $\intint{-H}{H}$ berjumlah paling banyak $d\,(2H + 1)^{d+1}$.
        (Kehinggaan inilah yang membuat pertanyaan 17 berjalan.)
  \item Sintesis, satu kalimat untuk masing-masing: (i) tunjukkanlah satu-satunya
        bahan analisis pada bukti Liouville (yaitu teorema bab ini yang
        mana, yang dipakai di mana); (ii) nyatakanlah ketegangan
        yang menggerakkannya (karena kebulatan memaksa $\abs{P(p/q)} \geq
        q^{-d}$, sedangkan kemulusan melarang $\abs{P(p/q)} > M\abs{x -
        p/q}$); (iii) bandingkanlah bukti Liouville dan bukti Cantor
        atas keberadaan \omterm{pb:b1:logic:1}{bilangan transenden}; (iv) sebutkanlah
        tempat jilid ini bertemu temanya lagi --- yaitu soal akhir
        pekan \cref{ch:b1:integration} yang membuktikan bahwa $\pi$
        irasional lewat apitan kebulatan lawan kekecilan yang sama,
        dengan integral menggantikan \omterm{def:b1:derivative:def}{turunan}.

\end{enumerate}
\end{problem}
